\documentclass[11pt]{article}
\usepackage[margin=0.85in]{geometry}
\usepackage{amsmath,amssymb,booktabs,longtable,xurl,hyperref}
\hypersetup{colorlinks=true,urlcolor=blue}
\setlength{\parindent}{0pt}
\setlength{\parskip}{0.55em}
\title{Fifty-Meter Perception: Full Scenario Campaign and Causal Failure Analysis}
\author{Collision Avoidance Research}
\date{October 3, 2026}
\begin{document}
\maketitle

\textbf{Main result.} All 14 requested noisy-sensor NRMM trials failed to return
the first solved control at $t=0$. No plant hold was executed. Zero recorded
collisions therefore provides no evidence of successful avoidance. The actual
conic solver reported primal infeasibility; no theoretical-assumption check
rejected the trials. A separate exact-state control experiment recovered in
11/14 cases and stopped without a solved input in three. All 694 MATLAB tests
passed. Unit-test success does not imply that these scenario experiments pass.

\section{Implemented experiment and range contract}
The sensor now publishes radar relative position precisely when
\[
 \|p_T-p_E\|_2\le 50\ \mathrm m.
\]
This is a reference-position radius, not rectangle clearance. It covers every
azimuth. Outside the disk, the sensor provides no target position and the
controller treats the target as absent. A complete current empty scan clears
its old forecast while preserving the previous input initialization. A generic
dropout without such a complete scan retains the previous forecast. Predicted
collision rows use the position radius at each node, including later re-entry;
the old permanent first-exit truncation was removed. Uncertain predictions are
active when their position enclosures can intersect the disk.

The actual estimator path in \path{runNonlinearPredictiveSafetyValidation.m}
uses \path{nrmmEstimatorControllerAdapter.m}. Controller inputs are observer
publications; the nonlinear plant always advances from its true state. Target
truth follows the original constant tangential acceleration and constant
sideslip trajectory, independently of each updated forecast. No synthetic
estimate perturbation stands in for the NRMM in the main experiment.

Seven scenarios (head-on, accelerating head-on, braking lead, crossing,
turning crossing, curved head-on, curved crossing) were tested at 8 and
15 m/s. They all start separated and collide under the independent given-path,
constant-speed baseline. Initial target ranges are 6.25--36.75 m: all 14 main
trials have a real in-range radar observation at the first controller call.
Consequently the common startup failure is not late detection at 50 m.

The controller period is 50 ms, the PCBF prefixes are 8 and 16 holds, and
completion adds 4 s, producing 4.4 and 4.8 s predictions. The 0.20 m sampled
collision margin, 10 m center-deviation limit, single CLF and single controller
are unchanged. There are no road-boundary constraints or executable backup
controls. Estimated $A$ and $\beta$ are fixed within each prediction and can
change at the next observation. No assumption-status diagnostics were added.

The estimator uses 80 Hz samples, at most 2.5 ms RK4 substeps, 0.5 s of causal
sensor history, and seed 20261003 for each case. The existing bounded-uniform
sensor maxima are GNSS position 0.04 m, GNSS velocity 0.05 m/s, IMU acceleration
0.03 m/s$^2$, gyro 0.0015 rad/s, and radar position 0.04 m (vector maxima are
Euclidean radii). The ego rear-axle parameter is copied from the 1.65 m plant.
Observer gains, operating-domain bounds, and the common 15 m/s target-speed
prior were not retuned to individual truth trajectories.

Experiments stop on collision, inability to issue a solved control, or observed
nominal recovery. Recovery requires a 1 s dwell with errors below
$(0.1\,\mathrm m,1^\circ,0.1\,\mathrm{m/s},0.05\,\mathrm{m/s},0.01\,\mathrm{rad/s})$.
The 2000-hold cap is only an observation limit, never a success requirement;
no case terminated merely at that cap. Recovery must follow the last nominal
encounter in that observation window. Circular constant-$\beta$ scenarios can
meet again: the long curved-head-on controls include a second encounter.
The replay checks 31 points per hold using \path{ode45} with tolerances
$10^{-11}$ and $10^{-12}$. These are sampled experiment results, not an
infinite-horizon or continuous-time certificate.

\section{Two independent contradictions in the main experiment}
\subsection{The uncertain endpoint cannot fit inside the nominal core}
The actual endpoint cone generated by \path{solvePredictiveControl.m} is
\[
 \|L_f(e_M+J_Md)\|_2\le r_f-R_M,
 \qquad R_M=\sum_j\|L_fG_{M,j}\|_2.
\]
Its right side has no decision-dependent term. In every case $R_M>r_f$, so
its right side is negative and its left side is nonnegative for every input.
This is an explicit infeasibility proof for the assembled problem, independent
of solver tuning, iteration count, initialization direction, or target motion.

\begin{center}
\begin{tabular}{rrrr}
\toprule
Ego speed & Core radius $r_f$ & Required $R_M$ & $R_M/r_f$\\
\midrule
8 m/s & $2.4414\times10^{-4}$ & 0.7316--0.8839 & 2997--3620\\
15 m/s & $9.7656\times10^{-4}$ & 0.4284--0.5838 & 439--598\\
\bottomrule
\end{tabular}
\end{center}
These are weighted-core norm units, not meters. The core's pure longitudinal
speed half-width is only $3.2172\times10^{-5}$ m/s at 8 m/s and
$1.2872\times10^{-4}$ m/s at 15 m/s. Published initial longitudinal-speed
radii are about 0.057--0.067 m/s. The reported lateral-speed radius is
0.46834 m/s, while the actual straight-road lateral-speed estimate error is
only 0.00167 m/s for this seed. The mathematical tightening concerns the full
bound, not the error that happened to be realized.

The design mismatch is precise: the implementation propagates ego generators
with open-loop $G_{i+1}=A_iG_i$, then asks that whole set to enter the tiny core
constructed for nominal state-feedback contraction. It does not propagate an
output-feedback error tube compatible with that terminal policy and its
nonzero measurement uncertainty. Removing the target entirely in an offline
same-frame counterfactual still produced no solved input in all 14 cases.
Thus a target-free controller already encounters the endpoint contradiction.

\subsection{The target prediction discards available directional information}
Setting only the ego error radius to zero also left all 14 problems infeasible.
The independent target problem is visible in the published prediction set:
initial course radius is $\pi$, speed interval is $[5,20]$ m/s, acceleration
interval is $[-2,2]$ m/s$^2$, and sideslip domain is $\pm0.005$ rad.
The current isotropic propagation uses
\[
 R_p(\tau)=R_{p,0}+R_V\tau+\tfrac12 R_A\tau^2
     +|s(\tau)|R_\gamma+\tfrac12 s(\tau)^2R_\kappa,
 \quad s(\tau)=\widehat V\tau+\tfrac12\widehat A\tau^2.
\]
Here $\widehat V=15$, $\widehat A=0$, $R_V=10$, $R_A=2$, and
$R_\gamma=\pi$. The course term alone contributes 207.35 m at 4.4 s and
226.19 m at 4.8 s. Total final position radii are 277.55--305.67 m; collision
row tightening reaches approximately 283--312 m. This is not a change of
target motion assumption: every member still has constant $A$ and $\beta$.
It is excessive uncertainty about those trajectories and their initial course.

The information loss occurs in \path{estimator/nrmmControllerErrorBounds.m}.
The function first intersects component and heading bounds with
\path{measurementHistoryEnclosure}, then constructs
\path{predictionErrorSet} from the older coarse component radii.
All 14 snapshots contain 41 history samples with an available heading interval,
but the prediction set still receives a whole-circle course radius. For the
15 m/s braking-lead case, the history heading radius is 0.09755 rad and the
published heading-component bound is 0.10572 rad; the prediction still uses
$\pi$. The intersection must be carried into the relative prediction set with
the appropriate frame correlations, not discarded or copied without a frame
transformation argument.

The target-only counterfactuals also provide direct row contradictions inside
the finite correction box. For example, in the 15 m/s head-on problem one hard
row demands $a^Td\le-201.16855$, whereas its minimum over that entire box is
$-2.67200$. The constraints therefore cannot be repaired by the current local
step. The contradiction remains even after the ego terminal reserve vanishes.

The declared target sideslip domain is also narrower than the turning-crossing
and curved-crossing truth values (0.05 and 0.04 rad). This does not reject any
trial and is not used as an experimental failure criterion. It is a separate
coverage limitation to address before interpreting these prediction bounds as
physical robust guarantees; it is not the cause of the already-proved empty
endpoint cones.

\section{Initialization further distorts the predicted encounter}
Even with 41 position observations, the adapter fits only the velocity
\emph{direction} and then normalizes its magnitude to the configured 15 m/s
prior. It initializes acceleration to zero. Consequently all 14 first
publications have speed 15 m/s and $A=\beta=0$, despite targets having different
speeds, acceleration and turning motion. The branch in
\path{localTargetStateFromInertialWindow} explicitly overwrites the magnitude
of the least-squares velocity. This is not an unavoidable consequence of
position-only sensing: a multi-sample window already supplies speed information.

Removing all uncertainty tightening while retaining those estimated point
states yields a solved first input in only 7/14 offline probes. This is not a
closed-loop success rate and these inputs were never issued. It shows why
removing or enlarging the terminal cone alone does not establish that the whole
observer/controller integration will work. The startup estimate and RTI
formulation also need attention.

\section{Exact-state control experiment}
The following experiment uses exact point states with the same 50 m visibility
rule. It is a separate comparison, not a substitute for the requested observer
experiment. There were no sampled collisions in the executed holds. Eleven
cases recovered; three stopped without a solved input. Minimum observed gaps
below 0.20 m remain possible because that value is an affine sampled constraint,
not a certified nonlinear or intersample gap.

\begin{center}\small
\begin{tabular}{rlrrl}
\toprule
Speed & Scenario & Stop time (s) & Min gap (m) & Outcome\\
\midrule
8 & Head-on & 10.35 & 1.3972 & Recovered \\
8 & Accelerating head-on & 10.30 & 1.0645 & Recovered \\
8 & Braking lead & 14.65 & 0.5526 & Recovered \\
8 & Crossing & 10.55 & 0.1482 & Recovered \\
8 & Turning crossing & 13.35 & 0.1936 & Recovered \\
8 & Curved head-on & 88.45 & 1.7952 & Recovered \\
8 & Curved crossing & 9.25 & 0.1931 & No solved input \\
15 & Head-on & 7.40 & 0.2345 & Recovered \\
15 & Accelerating head-on & 0.25 & 13.4433 & No solved input \\
15 & Braking lead & 5.25 & 1.2981 & No solved input \\
15 & Crossing & 8.20 & 1.3919 & Recovered \\
15 & Turning crossing & 8.90 & 1.3681 & Recovered \\
15 & Curved head-on & 63.95 & 2.8642 & Recovered \\
15 & Curved crossing & 8.20 & 1.0687 & Recovered \\
\bottomrule
\end{tabular}
\end{center}

\subsection{Accelerating head-on, 15 m/s, at 0.25 s}
Current rectangle clearance is 13.4433 m. The inherited-budget CLF solve returns
an affine solution, but its pose prediction error is 10.696 mm against a 10 mm
limit. Damping cannot start at the shifted anchor: its terminal-cone residual
is 0.031651. The replacement potential-field model has a feasible affine PCBF
solution, but the full CLF step has 36.471 mm error. Three damped trials reduce
this to 21.175, 18.662 and 18.007 mm, all above the limit.

The primary endpoint of this interpolation itself has 17.756 mm nonlinear
pose disagreement. Taking the secondary fraction to zero approaches this
nonlinear-inaccurate primary point, not the original nonlinear anchor.
Therefore an $O(\alpha^2)$ error-to-zero claim does not apply to this segment.

\subsection{Braking lead, 15 m/s, at 5.25 s}
Current rectangle clearance is 2.6789 m. The shifted CLF solve reaches zero CLF
slack, but its 10.124 mm prediction error exceeds the 10 mm limit. The shifted
base is unavailable for damping because its terminal residual is 0.017754.
The fresh full step has 14.095 mm error. Damping makes it worse: 20.948,
26.118 and 29.180 mm, approaching a primary-point error of 32.608 mm. This is
another geometrically misplaced interpolation base, not a conic-solver timeout.

\subsection{Curved crossing, 8 m/s, at 9.25 s}
Current target range is 39.5844 m and current rectangle clearance is 36.8107 m.
The shifted convex solution reaches essentially zero CLF slack but has
19.543 mm prediction error. Its damping base again violates the terminal cone
(residual 0.015317). The fresh potential-field seed predicts body overlap
between absolute times 12.575 and 12.750 s in the hard completion tail.
Ordinary configuration distance and its selected dual derivative are zero at
that overlap. A generated collision row is exactly
\[
 0\cdot d\le-0.20.
\]
No PCBF prefix slack appears in that tail row. Increasing the input correction
box cannot change its zero coefficients, and both fresh PCBF attempts report
primal infeasibility. Failure of this local convexification does not prove
that no physically safe maneuver exists.

The shared exact-state defect is that a shifted affine result is not necessarily
a feasible nonlinear continuation, particularly at a very small terminal core.
Once its terminal membership is lost, the inherited feasible segment used for
cheap damping disappears. The replacement primary can itself be nonlinearly
inaccurate or have a flat collision linearization. A second numerical solve or
smaller secondary fraction does not automatically resolve either case.

\section{Design implications and scope of this change}
The identified priorities are (i) an uncertainty propagation and terminal set
that describe the same closed-loop error dynamics, (ii) a measurement-consistent
target parameter set that preserves the history information and its correlations,
and (iii) initialization that retains measured speed information. These address
the two independent empty feasible sets and the incorrect initial forecast.
They do not require a second controller, a second CLF, or an executable backup.

The RTI correction also needs a nonlinear-consistent interpolation base and a
collision model that can restore an overlapping search seed. Merely deleting
the 10 mm test, setting uncertainty to zero, increasing the iteration budget,
or declaring missing theory assumptions to be failures would not repair these
structural issues. Those counterfactuals are not implemented as new runtime modes.
This task implements the requested sensing contract and produces evidence;
it does not claim that the failure mechanisms have already been fixed.

\section{Validation and reproducibility}
MATLAB ran \path{results=runtests('tests'); assertSuccess(results)} using
\path{matlab -singleCompThread -batch}. The suite completed with 694 passed, zero failed and zero incomplete tests. The test-runner
duration was 282.043912 s; this is not user work time. New cases verify the
50 m boundary in three azimuths, no outside-disk position publication, clearing
an old target after a complete empty scan while retaining the input warm start,
reference-position versus body-clearance radius, and later predicted entry.
Existing generic dropout and recovery/failure semantics still pass.

MATLAB Code Analyzer was run on the eight relevant source/test files. It
reported no new code issues; existing sparse-indexing notices remain. The local
Analyzer preference file was unavailable, so Analyzer used its default settings.
Git whitespace checks passed. Runtime comparisons were not isolated from other
active processes and do not establish a real-time deadline.

The campaign entry point is \path{scripts/runPotentialFieldCampaign.m} with
\path{EstimatorConfiguration=cfg}, where
\path{cfg=estimatorControllerIntegrationConfig()} and
\path{cfg.randomSeed=20261003}. The exact-state control omits that option.
Both use 2000 as an observation cap, not a required duration. The actual main
run consumes noisy sensor samples, not truth kinematics. The separate script
\path{scripts/analyzePerceptionRangeFailures.m} reconstructs failed calls and
uses non-stopping conditional breakpoints in a dedicated MATLAB process to
capture actual conic matrices and candidates. Its target-free, ego-radius-zero
and point-estimate-only probes are offline ablations, never executable controls.

Compact tables, full causal-frame records, settings and test results are in
\path{report/PERCEPTION_50M_CAMPAIGN_20261003/}. Original MAT snapshots,
continuations, full JSON trajectories, commands and logs are retained in the
identified external experiment export. No figure, external solver dependency,
generated native binary, unrelated estimator report or administrative file is
included in this task's project commit.
\end{document}
